\begin{figure*}[p]
\centering
% SINGLE DOWNWARD FLOW ON A FIXED GRID.
%
% THE RECURRING FAILURE HERE, THREE TIMES NOW: minimum height is a FLOOR, not a height.
% A box whose text wraps to one more line than expected grows downward into the row beneath
% and prints through it, and this is invisible in the source. The rule that finally works:
% EVERY BOX IS AT MOST TWO LINES. Text is kept short enough to guarantee that at the box
% width given, anything longer goes in the caption, and rows are then spaced at twice the
% two-line height. Detail hangs to the RIGHT of the version boxes where there is free width,
% not inside them where it would add lines.
%
% Also fixed:
%   * v3|v4|v5 ran ACROSS, so the arrow into v3 came up from below-left. A version history
%     is a sequence and now reads as one.
%   * the previous attempt was drawn ~200mm wide and shrunk by \maxwidth to fit, which is
%     why the type came out at 6pt in a field of white space. The design width is ~175mm
%     against a ~165mm \linewidth, so \maxwidth scales by about 0.94 and the type stays
%     essentially the size it was set in.
%   * the stub joining v5 to its description was drawn THROUGH the description.
%   * jargon: "geometric invariants" told the reader nothing. The gates now say what is
%     checked, in words a surgeon would use.
%   * "out" is a reserved TikZ key: naming a node (out) silently drops its shape and
%     surfaces much later as "missing \item". The final node is (rel).
\maxwidth{%
\begin{tikzpicture}[
  x=1mm, y=1mm,
  font=\small,
  every node/.style={inner sep=0pt},
  bx/.style={draw, rounded corners=1.5pt, align=center, text width=#1,
             minimum height=11mm, inner sep=2.5pt},
  src/.style={bx=52mm, fill=black!5},
  proc/.style={bx=168mm, fill=blue!7},
  rec/.style={bx=38mm, fill=black!3},
  ver/.style={bx=24mm, fill=green!14, minimum height=9mm},
  det/.style={align=left, text width=124mm, font=\small, anchor=west},
  gate/.style={bx=168mm, fill=orange!14},
  fin/.style={bx=168mm, fill=green!22},
  ar/.style={-{Stealth[length=5pt]}, semithick, black!70},
  ln/.style={semithick, black!70},
]

% --- row 1: sources ------------------------------------------- 2 lines, half 7
\node[src] (tcia)   at (-58, 0) {\textbf{TCIA CT Colonography}\\[2pt]the scans};
\node[src] (spine)  at (  0, 0) {\textbf{CTSpine1K}\\[2pt]vertebra outlines};
\node[src] (pelvic) at ( 58, 0) {\textbf{CTPelvic1K}\\[2pt]sacrum + hip outlines};
\draw[ln] (tcia.south) -- (-58,-15) -- (0,-15);
\draw[ln] (pelvic.south) -- (58,-15) -- (0,-15);
\draw[ln] (spine.south) -- (0,-15);
\draw[ar] (0,-15) -- (0,-22);

% --- row 2: the crosswalk ------------------------------------- centre -30, half 8
\node[proc] (cross) at (0,-30) {\textbf{match each set of outlines to the scan it was drawn
  on}\\[2pt]every patient was scanned twice and neither source recorded which scan it used};

% --- row 3: the four kinds of record -------------------------- centre -56, half 9
\node[rec] (fused) at (-60,-56) {\textbf{fused}\\[2pt]both, one scan\\[1pt]$n{=}342$};
\node[rec] (sep)   at (-20,-56) {\textbf{separate}\\[2pt]pelvis on the other
  scan\\[1pt]$n{=}351$};
\node[rec] (sonly) at ( 20,-56) {\textbf{spine only}\\[2pt]no pelvis
  outline\\[1pt]$n{=}89$};
\node[rec] (ponly) at ( 60,-56) {\textbf{pelvis only}\\[2pt]held back\\[1pt]$n{=}20$};
\draw[ln] (cross.south) -- (0,-44);
\draw[ln] (-60,-44) -- (60,-44);
\foreach \x in {-60,-20,20,60} \draw[ar] (\x,-44) -- (\x,-47);

% --- row 4: corrections --------------------------------------- centre -82, half 8
\node[proc] (fix) at (0,-82) {\textbf{fix mistakes in the original
  outlines}\\[2pt]one vertebra split in two, a level skipped, a level left unlabelled:
  $n{=}112$ spine fixes over 103 patients, $n{=}13$ pelvis fixes over 12};
\foreach \x in {-60,-20,20,60} \draw[ln] (\x,-65) -- (\x,-70);
\draw[ln] (-60,-70) -- (60,-70);
\draw[ar] (0,-70) -- (0,-74);

% --- row 5: filling in the missing half ----------------------- centre -108, half 8
\node[proc] (dense) at (0,-108) {\textbf{fill in whichever half is
  missing}\\[2pt]a pelvis has nothing to count so a model can fill it in; the 20
  pelvis-only spines were filled in and then corrected \emph{by hand}};
\draw[ar] (cross.south -| 0,0) ++(0,0) -- (0,-90) -- (0,-100);

% --- rows 6-8: the releases, in order -------------------------- 20mm apart
\node[ver] (v3) at (-72,-132) {\textbf{v3}};
\node[ver] (v4) at (-72,-152) {\textbf{v4}};
\node[ver] (v5) at (-72,-172) {\textbf{v5}};
\draw[ar] (0,-116) -- (0,-122) -- (-72,-122) -- (v3.north);
\draw[ar] (v3.south) -- (v4.north);
\draw[ar] (v4.south) -- (v5.north);

\foreach \y in {-132,-152,-172} \draw[ln, black!35] (-60,\y) -- (-54,\y);
\node[det] at (-52,-132) {added the thigh bones and the lower chest vertebrae, and split
  the top block of the sacrum (S1) off as its own label};
\node[det] at (-52,-152) {added the \textbf{ribs}, merged from two automatic tools, then
  sent $n{=}152$ scans to a person to check and correct};
\node[det] at (-52,-172) {gave a \textbf{rib growing from a lumbar vertebra} its own label
  instead of calling it the twelfth rib ($n{=}15$ scans)};

% --- row 9: the checks ----------------------------------------- centre -196, half 8
\node[gate] (qc) at (0,-196) {\textbf{checks every scan must pass, in order}\\[2pt]outlines
  sit on the bone and left is on the left $\rightarrow$ each rib touches the vertebra its
  number says $\rightarrow$ measurements fall in human ranges};
\draw[ar] (v5.south) -- (-72,-186) -- (0,-186) -- (qc.north);

% --- row 10: release ------------------------------------------- centre -220
\node[fin] (rel) at (0,-220) {\textbf{802 released scans}\\[2pt]spine, pelvis, ribs and
  thigh bones on one image; 33 carry a radiologist's Castellvi grade};
\draw[ar] (qc.south) -- (rel.north);
\end{tikzpicture}}%
\caption{How the dataset was built. Three public collections are combined patient by
patient. Neither annotation source recorded which scan its outlines were drawn on, and every
patient in this cohort was scanned twice, so each set of outlines has to be matched to the
scan whose bone it actually lands on --- an outline placed on the wrong scan of the right
patient is a correct segmentation that simply does not line up with the image beneath it.
The four kinds of record are the result of that matching and are the main finding of it: in
351 patients the two collections had drawn their outlines on \emph{different scans of the
same person}. Numbers marked $n$ are scans, or corrections, at that step.}
\label{fig:pipeline}
\end{figure*}
